\documentclass[11pt]{article}

\usepackage[utf8]{inputenc}
\usepackage[T1]{fontenc}
\usepackage{amsmath, amssymb}
\usepackage{graphicx}
\usepackage{booktabs}
\usepackage[margin=1in]{geometry}
\usepackage[hidelinks]{hyperref}
\usepackage{caption}
\usepackage{fancyhdr}

\graphicspath{{outputs/plots/}}

% --- running course header, like the course template ---------------------- %
\pagestyle{fancy}
\fancyhf{}
\lhead{\small Data Driven Modelling of Dynamical Systems}
\rhead{\small \href{https://github.com/lfpc/ode_modelling}{\texttt{github.com/lfpc/ode\_modelling}}}
\cfoot{\thepage}
\renewcommand{\headrulewidth}{0.4pt}

\newcommand{\dEdx}{\langle dE/dx\rangle}
\newcommand{\unit}[1]{\,\mathrm{#1}}

\title{\vspace{-2em}Modelling Muon Transport in Iron\\ with Neural ODEs and SDEs}
\author{Luis Felipe Cattelan \\[2pt]
        \small Physik-Institut, University of Z\"urich \\[1pt]
        \small Course project --- Data Driven Modelling of Dynamical Systems}
\date{\today}

\begin{document}
\maketitle
\thispagestyle{fancy}

\section{Introduction}
When a charged particle crosses a block of matter it does not travel in a
straight line at constant energy: it continuously loses momentum to the atomic
electrons and is randomly deflected by the nuclei. Detector simulations such as
Geant4 reproduce this in great detail, collision by collision, but at a high
computational cost. The idea of this project is to ask whether a small neural
network, trained on Geant4 output, can act as a fast \emph{surrogate} for one
elementary process: the change in a muon's momentum after it travels a short
step through iron. Practically, such a learned per-step map could replace the
expensive Monte-Carlo stepping inside a simulation; more importantly, it is a clean test for a conceptual question about deterministic
versus stochastic modelling.

The model I want to learn is best stated by looking at the physics. A muon of
momentum $p$ that travels an arc length ``$\mathrm{step}$'' through iron undergoes
two things:
\begin{itemize}
  \item it \textbf{loses momentum along its direction of travel} ($dP_z<0$) by
  ionizing the atomic electrons
  \item it is \textbf{deflected sideways} ($dP_t$) by many small-angle Coulomb
  scatters off the nuclei
\end{itemize}
The object to be learned is therefore the map
$(p_0,\,\mathrm{step})\mapsto p_1$, the momentum after one step. The catch is that
the data are stochastic -- every muon scatters differently -- while a
deterministic network outputs only one value per input. Can such a model still
learn the process ``in the mean'', and where does it break down?

More formally, the muon's momentum along its path obeys an equation of motion of
the form
\begin{equation}
  \frac{dp}{ds} \;=\;
  \underbrace{q\,\hat{p}\times B}_{\text{Lorentz force (known)}}
  \;+\;
  \underbrace{f_\theta(p,\,\mathrm{step})}_{\text{material transport (to learn)}},
  \label{eq:eom}
\end{equation}
where $\hat p=p/|p|$: the first term is the deflection by a magnetic field $B$ and
the second is the change caused by the material (energy loss and scattering). The
Lorentz term is known analytically and trivial to evaluate once the field is given, so the genuinely unknown, hard
part is the material transport. The goal of this project is precisely to learn
that transport term, $f_\theta$, from the Geant4 data.

\section{The model}
\subsection{Overview}
I build two models that share the same neural backbone but differ in their output
and training objective: a \emph{deterministic} Neural ODE that learns the mean
transport rate $dp/ds$, and a \emph{stochastic} Neural SDE that learns a drift
and a diffusion. Both are small multilayer perceptrons (three hidden layers,
$64$ units, $\tanh$) implemented in PyTorch. Before describing them I recall the
reference physics they are compared against.

\subsection{Reference physics (the formulae we compare to)}
\paragraph{Mean energy loss -- Bethe--Bloch \cite{pdg}.}
\begin{equation}
  -\left\langle\frac{dE}{dx}\right\rangle
   = K z^2 \frac{Z}{A}\frac{1}{\beta^2}
     \left[\tfrac12\ln\frac{2 m_e c^2\beta^2\gamma^2 W_{\max}}{I^2}
           - \beta^2 - \frac{\delta(\beta\gamma)}{2}\right],
  \label{eq:bb}
\end{equation}
with $K=0.307\unit{MeV\,cm^2/mol}$ and, for iron, $Z/A=0.466$,
$\rho=7.87\unit{g/cm^3}$, $I=286\unit{eV}$; $W_{\max}$ is the maximum
single-collision energy transfer and $\delta$ the Sternheimer density-effect
correction. For a relativistic muon $dE\approx c\,dP_z$, so
$\langle|dP_z|\rangle\propto\mathrm{step}$.

\paragraph{Multiple scattering -- Highland \cite{highland}.}
\begin{equation}
  \theta_0 = \frac{13.6\unit{MeV}}{\beta c\,p}\,z\,\sqrt{x/X_0}
             \left[1+0.038\ln\!\big(x z^2/(X_0\beta^2)\big)\right],
  \label{eq:highland}
\end{equation}
with $X_0=1.76\unit{cm}$ for iron. The transverse momentum kick is
$p_T\approx p\,\theta$, so $\mathrm{RMS}\,|p_T|=\sqrt2\,p\,\theta_0\propto\sqrt{x}$.

\paragraph{Energy-loss fluctuations -- Landau/Vavilov \cite{landau}.}
The energy loss in a thin layer is Landau-distributed (a sharp peak with a long
tail from rare hard $\delta$-rays). Its variance (Bohr) is
\begin{equation}
  \sigma_E^2 = \tfrac{K}{2}\,\frac{Z}{A}\,\frac{z^2}{\beta^2}\,\rho x\,
               W_{\max}\,\big(1-\tfrac{\beta^2}{2}\big).
  \label{eq:vavilov}
\end{equation}

\subsection{Data and simulation setup}
The data were produced with Geant4 \cite{geant4}, which transports muons through
iron and stores, step by step along each trajectory, the momentum before and
after the step; every row of the resulting table is one such single step. I work
in a deliberately reduced setting: a single material (iron), no magnetic field
(a field-free slab, so the only momentum changes come from the material) and no
tracking of the spatial position -- only the momenta. Of the physics I keep
ionization energy loss ($dP_z$) and multiple Coulomb scattering ($dP_t$);
radiative losses (bremsstrahlung etc.) are $\lesssim$ a few percent at these
energies and are not modelled separately.

I take the muon to start along $z$, so the initial momentum is
$p_0=[0,0,p_0]$ with $p_0\in[10,20]\unit{GeV}$, and after one step
$p_1=[dP_t,\,0,\,p_0+dP_z]$: the longitudinal component shrinks by the energy
loss and a transverse kick of magnitude $dP_t$ appears. The kick is placed along
$x$ because its azimuth is physically uniform and carries no information; the
random direction is restored at inference by an azimuthal rotation. Concretely we
simulate millions of muons: for each one a step size is drawn, and Geant4 returns
that (random) step together with the resulting $dP_z$ and $dP_t$.

For training we feed the network the logarithm of the step size and normalize all
inputs, which is necessary for numerical stability.

%\begin{figure}[t]
%  \centering
%  \includegraphics[width=\linewidth]{data_overview.png}
%  \caption{Exploratory analysis of the dataset. Top: step-size and momentum
%  distributions, and the mean relative change vs.\ step (log--log) with reference
%  slopes $1$ (Bethe--Bloch) and $\tfrac12$ (scattering). Bottom: the energy-loss
%  rate and the scattering rate, normalized so that each is flat if the
%  corresponding power law holds.}
%  \label{fig:data}
%\end{figure}

\subsection{Deterministic Neural ODE}
The motion over arc length $s$ obeys an equation of the form $dp/ds=f(p,s)$.
Instead of hand-coding $f$ from Bethe--Bloch and scattering theory, a Neural ODE
\cite{neuralode} \emph{learns} it: the rate is parametrized by the network,
\begin{equation}
  \frac{dp}{ds}=f_\theta(p,\,\mathrm{step}),
  \label{eq:node}
\end{equation}
and the momentum after a finite step is obtained by a fourth-order Runge--Kutta
(RK4) integration,
\begin{align}
  k_1 &= f_\theta(p_0,\,ds), \qquad
  k_2 = f_\theta\!\big(p_0+\tfrac{ds}{2}k_1,\,ds\big), \notag\\
  k_3 &= f_\theta\!\big(p_0+\tfrac{ds}{2}k_2,\,ds\big), \qquad
  k_4 = f_\theta\!\big(p_0+ds\,k_3,\,ds\big), \notag\\
  \hat p_1 &= p_0 + \tfrac{ds}{6}\big(k_1+2k_2+2k_3+k_4\big).
  \label{eq:rk4}
\end{align}
Training is supervised on the single-step transitions $(p_0,\mathrm{step},p_1)$:
for each sample we integrate one RK4 step and minimize the mean squared error of
the \emph{relative} momentum change,
\begin{equation}
  \mathcal{L}_{\mathrm{MSE}}
  = \Big\langle\, \big\|\,(\hat p_1-p_0)/|p_0| - (p_1-p_0)/|p_0|\,\big\|^2 \,\Big\rangle .
  \label{eq:mse}
\end{equation}
Dividing by $|p_0|$ makes the target dimensionless and of order $10^{-3}$
independent of the momentum. The key property is that the minimizer of an MSE
loss is the \emph{conditional mean}: $f_\theta\to\mathbb{E}[dp\mid p_0,
\mathrm{step}]$. The deterministic model therefore learns the mean energy loss
and the mean transverse kick, but discards all fluctuation around them. At
inference a trajectory is generated by chaining RK4 steps, with an optional
random azimuthal rotation after each step to mimic the random scattering
direction.

\subsection{Stochastic Neural SDE}
I model the same transport as an It\^o stochastic differential equation,
\begin{equation}
  dp = \mu_\theta(p,s)\,ds + \sigma_\theta(p,s)\,dW,\qquad dW\sim\mathcal N(0,ds\,\mathbb I),
\end{equation}
where the network outputs a drift $\mu$ and a (log-)diffusion $\sigma$. A single
Euler--Maruyama step reads $\Delta p=\mu\,ds+\sigma\sqrt{ds}\,\xi$ with
$\xi\sim\mathcal N(0,\mathbb I)$. The one-step transition is the diagonal Gaussian
$p_1\mid p_0\sim\mathcal N(p_0+\mu\,ds,\ \mathrm{diag}(\sigma^2 ds))$, and I train
by minimizing its negative log-likelihood (NLL). Unlike MSE, the NLL fits the
mean \emph{and} the spread. Because scattering is isotropic but the data fix the
kick along one axis, we apply a random azimuthal rotation to each $(p_0,p_1)$ pair
during training; this forces the transverse drift to vanish and makes
$\sigma_x=\sigma_y$, so the transverse spread is carried entirely by the
diffusion term.

\paragraph{Why an SDE is the natural object.} An ODE integrated by RK4 produces
increments $\Delta p\propto ds$. Scattering requires $\Delta p_t\propto\sqrt{ds}$,
i.e.\ a rate $dp_t/ds\propto1/\sqrt{ds}$ that diverges as $ds\to0$ -- impossible
for a bounded network. The SDE puts the $\sqrt{ds}$ in the noise term, so the
network only learns a smooth $\sigma$.

\subsection{Training, tools and sources}
Both models are trained on all $2\times10^6$ steps (a $90/10$ train/validation
split) with Adam (learning rate $10^{-3}$), batches of $1024$, gradient-norm
clipping at $1$, and a learning-rate-on-plateau schedule, for $10$--$15$ epochs;
I keep the lowest-validation checkpoint. I found the weight decay had to be set
to zero: at the usual $10^{-3}$ it is comparable to the signal itself (targets
$\sim10^{-3}$) and simply drives the predictions toward zero. The code is written
in Python with PyTorch \cite{pytorch} and is available at \url{https://github.com/lfpc/ode_modelling}; the analytic reference curves use the PDG
parametrizations \cite{pdg, highland}. 

\section{Analysis and validation}
\subsection{Mean energy loss vs.\ Bethe--Bloch}
Figure~\ref{fig:bb} compares the model's predicted mean $|dP_z|$ with the data
and Bethe--Bloch. The Neural ODE overlays the data across all step sizes. The
naive stopping power averaged over a step window sits $\sim25\%$ above the
Bethe--Bloch value of $15.5\unit{MeV/cm}$ -- but, as explained next, this is an
artifact of the data, not of the model: restricting to large, self-averaging
steps removes it. In the window $1$--$5\unit{cm}$ (Fig.~\ref{fig:mom}) the data
($15.47$), model ($15.25$) and theory ($15.48\unit{MeV/cm}$) agree at the percent
level and show the same relativistic rise with momentum.

\begin{figure}[t]
  \centering
  \includegraphics[width=\linewidth]{bethe_bloch_comparison.png}
  \caption{Energy loss vs.\ Bethe--Bloch. Left: mean $|dP_z|$ vs.\ step. Centre:
  stopping power vs.\ step (the small-step blow-up is the selection bias of
  Sec.~\ref{sec:bias}). Right: stopping power vs.\ momentum in a step window.
  The NeuralODE (orange) tracks the data (blue) everywhere.}
  \label{fig:bb}
\end{figure}

\begin{figure}[t]
  \centering
  \includegraphics[width=0.6\linewidth]{bethe_bloch_vs_momentum_largestep.png}
  \caption{Stopping power vs.\ momentum, restricted to the self-averaging window
  ($1$--$5\unit{cm}$): data, model and Bethe--Bloch agree to $\sim1\%$.}
  \label{fig:mom}
\end{figure}

\subsection{The step-length selection bias}
\label{sec:bias}
The stopping power is a \emph{mean over many collisions}; a step recovers it only
if it contains enough of them. For short steps the loss is Landau-distributed and
dominated by whether a rare hard collision occurred. Moreover, Geant4 places step
boundaries \emph{at} interaction points, so short steps are preferentially those
containing a hard $\delta$-ray (biased high) while long steps are
soft-collision-only stretches (biased low). Measuring the local
$\langle|dP_z|\rangle/\mathrm{step}$ in narrow windows confirms this: it falls
from $\sim500\unit{MeV/cm}$ at $0.1\unit{mm}$, through the Bethe--Bloch value at
$1$--$3\unit{cm}$, down to $\sim9\unit{MeV/cm}$ at $10\unit{cm}$. The sub-linear
fit $dP_z\propto\mathrm{step}^{0.62}$ is the fingerprint of this artifact, not of
the physics. The model is faithful; the bias lives in the data, and it is removed
by evaluating only where the data self-averages.

\subsection{Scattering vs.\ Highland}
The left panel of Fig.~\ref{fig:diff} compares the SDE's learned diffusion with
the data RMS$|p_T|$ and Highland. In the self-averaging window the three agree
closely: data $23.3$, model $22.8$, Highland $21.5\unit{MeV}$. The $\sim8\%$
excess of the data over Highland is expected: Highland describes only the
Gaussian core, whereas the RMS is sensitive to the non-Gaussian large-angle
(Rutherford) tail. Removing the top $2\%$ of events (which carry $14.5\%$ of the
variance) brings the data core to $21.5\unit{MeV}$, matching Highland exactly.

\subsection{Energy-loss straggling vs.\ Vavilov}
The right panel of Fig.~\ref{fig:diff} compares the SDE's $\sigma_z\sqrt{ds}$ with
the data spread and Bohr/Vavilov, Eq.~\eqref{eq:vavilov}. The theory
($70.5\unit{MeV}$) is within $\sim13\%$ of the data ($81.3\unit{MeV}$) with the
correct $\sqrt{\mathrm{step}}$ slope, but the Gaussian SDE underestimates it
($53.4\unit{MeV}$). This is the expected failure of a Gaussian to represent a
heavy Landau tail: the likelihood-optimal Gaussian width is anchored to the bulk
and cannot account for the rare catastrophic losses.

\begin{figure}[t]
  \centering
  \includegraphics[width=\linewidth]{sde_diffusion_comparison.png}
  \caption{SDE diffusion vs.\ fluctuation theory. Left: transverse width vs.\
  Highland (excellent agreement). Right: longitudinal width vs.\ Bohr/Vavilov;
  the Gaussian SDE captures the scale but underestimates the heavy-tailed width.}
  \label{fig:diff}
\end{figure}


\section{Discussion}
This project started from a simple question -- can a deterministic network learn
a stochastic physical process ``in the mean'' -- and ended up being a clean
demonstration of the difference between drift and diffusion. The most satisfying
result is that, after fixing two non-obvious issues (the input normalization, and
restricting the comparison to self-averaging steps), the learned energy loss lands
essentially on top of the data from Geant4. 

The two main difficulties were as instructive as the successes. First, the
\emph{step-length selection bias} in the Geant4 data: I initially read the
$\sim25\%$ excess in the stopping power as a model error, and only understood it
after measuring the local $dE/dx$ as a function of step and realizing that short
steps in Geant4 are statistically special. A cleaner dataset (fixed step length,
or the continuous loss reported separately from $\delta$-rays) would avoid this
entirely. Second, the \emph{Gaussian-versus-Landau} limitation: the SDE captures
the variance of each channel but not the skew of the energy-loss distribution nor
the power-law scattering tail. A natural next step would be a heavy-tailed noise
(for example modelling $\log dP_z$, or using a Student-$t$ or normalizing-flow
noise), which I expect would close the remaining gap in the straggling.

A more principled framework, and the one I would adopt if I continued this work,
is the \emph{Neural Jump Stochastic Differential Equation} of Jia and Benson
\cite{njsde}. There the state evolves by a continuous neural flow (a drift, and
optionally a diffusion) but is also subject to discrete \emph{jumps} that occur at
random times drawn from a learned point process, each jump applying a learned
(stochastic) increment to the state. This matches the physics almost exactly: the
continuous flow is the soft, frequent collisions -- the smooth $dE/dx$ and the
Gaussian core of the scattering that my models already capture -- while the
\emph{jumps} are precisely the rare hard events I miss, namely the energetic
$\delta$-rays that build the Landau tail of the energy loss and the large-angle
Rutherford scatters that build the non-Gaussian tail of the deflection. Splitting
the dynamics into a smooth part and a jump part this way, rather than forcing a
single Gaussian to absorb both, is the natural way to reproduce the heavy tails
that are the one ingredient missing from the present model.

In summary, the mean energy loss is learned essentially exactly; the scattering
width is captured well by the Gaussian SDE; and only the heavy tail of the
energy-loss fluctuations is genuinely missed.

\begin{thebibliography}{9}
\bibitem{pdg} Particle Data Group, R.~L.~Workman \emph{et al.},
  ``Review of Particle Physics'' ,
  \emph{Prog. Theor. Exp. Phys.} (2022).
\bibitem{highland} V.~L.~Highland, ``Some practical remarks on multiple
  scattering,'' \emph{Nucl. Instrum. Meth.} \textbf{129}, 497 (1975);
  G.~R.~Lynch and O.~I.~Dahl, \emph{Nucl. Instrum. Meth.} \textbf{B58}, 6 (1991).
\bibitem{landau} L.~D.~Landau, ``On the energy loss of fast particles by
  ionization,'' \emph{J. Phys. USSR} \textbf{8}, 201 (1944).
\bibitem{geant4} S.~Agostinelli \emph{et al.},
  ``Geant4 -- a simulation toolkit,'' \emph{Nucl. Instrum. Meth.} \textbf{A506},
  250 (2003).
\bibitem{neuralode} R.~T.~Q.~Chen, Y.~Rubanova, J.~Bettencourt and D.~Duvenaud,
  ``Neural Ordinary Differential Equations,'' \emph{NeurIPS} (2018).
\bibitem{njsde} J.~Jia and A.~R.~Benson, ``Neural Jump Stochastic Differential
  Equations,'' \emph{NeurIPS} (2019).
\bibitem{pytorch} A.~Paszke \emph{et al.}, ``PyTorch: An Imperative Style,
  High-Performance Deep Learning Library,'' \emph{NeurIPS} (2019).
\end{thebibliography}

\end{document}
